\documentstyle[%


righttag%


,11pt%


,amssymb%


,russian%               remove in Eng.


,izvru%                 Set izven% in Eng.


]{amsart}





%





\newcommand{\la}{\langle}


\newcommand{\ra}{\rangle}


\newcommand{\tts}[1]{}





\theoremstyle{plain}


\theorembodyfont{\it}


\newtheorem{thms}{\hskip\parindent Теорема}[section]


\newtheorem{lems}{\hskip\parindent Лемма}[section]





\theoremstyle{definition}


\theorembodyfont{\rm}


\newtheorem{cornonum}{\hskip\parindent Следствие}


\newtheorem{examnonum}{\hskip\parindent Пример}


\newtheorem{notenonum}{\hskip\parindent Замечание}


\renewcommand{\thecornonum}{}


\renewcommand{\theexamnonum}{}


\renewcommand{\thenotenonum}{}





%\hoffset=-15mm%                  For Eng.


\setcounter{page}{55}


\god{2003}


\nomer{5~(492)} %                        Remove in Eng.


%\nomer{Vol.\,47, No.\,5}%               Set in Eng.


%\str{\pageref{begin}--\pageref{end}}%  Set in Eng.


\udk{517.98}





\author{\it М.Р.\,Тимербаев}


% NB:   ^^ remove in Eng.


\title{Пространства с нормой графика \\ и усиленные пространства


Соболева. I}


\thanks{Работа выполнена при финансовой


поддержке Российского фонда фундаментальных исследований (код


проекта 01-01-00616).}





\date{}


%\pagestyle{myheadings}%         Set in Eng.


\pagestyle{plain} %              Remove in Eng.


\begin{document}


%\thispagestyle{plain}%          Set in Eng.


\maketitle


\label{begin}








\section*{Введение}





Пусть $\Omega$ --- ограниченная плоская область с кусочно-гладкой


границей $\Gamma_0=\partial \Omega$, $\Gamma_i \subset \overline{\Omega}$,


$i=1,2,\dotsc,m$, --- гладкие кривые 


с концами на $\Gamma_0$. Известно, что если функция


$u$ принадлежит пространству Соболева $W_p^1(\Omega)$ ($1<p<\infty$),


то ее след $u|_{\Gamma_i}$ является элементом пространства


Соболева--Слободецкого $W_p^{1-1/p}(\Gamma_i)$. Пусть $V$ --- подмножество


функций из $W_p^1(\Omega)$, следы которых на $\Gamma_i$ обладают дополнительной


гладкостью, а именно, принадлежат, кроме того,


и классу $W_q^1(\Gamma_i)$ ($1<q<\infty$).


Определим на $V$ норму


\begin{equation}


{\left\| {u} \right\|}_V={\left\| {u} \right\|}_{W_p^1(\Omega)}+


\sum \limits_{i=0}^{m}


{\left\| {u} \right\|}_{W_q^1(\Gamma_i)}. \label{norma1}


\end{equation}


Далее, на цилиндре $Q=(0,1) \times \Omega$ определим множество функций


$W$ с конечной нормой


\begin{equation}


{\left\| {u} \right\|}_W = {\bigg({\int_{0}^{1} {\left\| {u(t)} \right\|}_


{_{W_p^1(\Omega)}}^{^{p}} dt} \bigg)}^{{1/p}} +


{\bigg( \sum_{i=0}^{m} {\int _{0}^{1} {\left\| {u(t)} \right\|}_


{_{W_q^1(\Gamma_i)}}^{^{q}} dt} \bigg)}^{{1/q}}. \label{norma2}


\end{equation}


Определенные выше пространства естественным образом возникают во многих


приложениях, к примеру,


при описании совместного движения поверхностных и грунтовых вод,


в гидродинамической теории смазки, 


в теории упругости \cite{Anton}--\cite{Daut}.


Вариационные задачи и задачи на собственные значения, а также их


конечноэлементные аппроксимации


в пространствах с нормой (\ref{norma1}) в случае многоугольной области


$\Omega$ и отрезков (стержней) $\Gamma_i$ при $p=q=2$ рассматривались


в \cite{Djak96}--\cite{DjakIzVuz}, где эти пространства


называются усиленными пространствами Соболева.





В данной статье пространства Соболева с усиленной метрикой


рассматриваются как частный случай следующей абстрактной


конструкции. Пусть имеются $B$-пространства $U$ и $Y$ и линейный


непрерывный оператор $\gamma \in L(U,Y)$. Для $B$-пространства $X


\subset Y$ определим пространство


$$


\gamma (U, X) \overset{\mathrm{def}}{=}{} \{u \in U: \gamma u \in X\},


$$


наделенное нормой графика


$$


{\|u\|}_{\gamma(U,X)}={\|u\|}_U+{\|\gamma u\|}_X.


$$


Многие пространства устроены подобным образом. В частности, для


пространств с нормами (\ref{norma1}) или (\ref{norma2}) роль


оператора $\gamma$ играет оператор следа на $\Gamma=\bigcup


\limits_{i=0}^{m} \Gamma_i$. В статье в предположении о


дополняемости ядра $\ker \gamma$ устанавливаются общие свойства


пространств $\gamma(U,X)$ в контексте свойств $U$ и $X$:


интерполяционное свойство, критерии компактности и плотности


множеств. Приводятся примеры реализаций описанной конструкции. В


пространствах с нормами (\ref{norma1}) и (\ref{norma2})


установлена плотность гладких функций.








\section{Предварительные результаты}





Через $\Phi$, $\Psi$ обозначаются топологические


векторные пространства (ТВП). Всюду, если не оговорено особо,


для ТВП включение $U \subset \Phi$ будет пониматься не


только в теоретико-мно\-жест\-вен\-ном, но и в топологическом смысле;


таким образом, равенство $U=V$ для $B$-прост\-ранств будет


означать также эквивалентность норм этих пространств.


Отметим также следующее: если $B$-пространства $U, V \subset


\Phi$ таковы, что $V$ является подмножеством $U$, то $V$


непрерывно вложено в $U$ --- это следует из теоремы Банаха о


замкнутом графике, которую нужно применить к тождественному


отображению.





Как обычно, через $L(\Phi, \Psi)$ обозначается множество


линейных непрерывных отображений из~$\Phi$ в $\Psi$. Пусть


$\gamma \in L(\Phi, \Psi)$ и $U \subset \Phi$. Тогда, очевидно,


$\gamma \in L(U, \Psi)$ (сужение $\gamma$ на $U$ здесь и далее


будет обозначаться тем же символом $\gamma$). Если $U \subset \Phi$ ---


$B$-пространство, то линейное множество $\gamma(U)$ (образ $U$


при отображении $\gamma$), наделенное фактор-нормой


$$


{\|x\|}_{\gamma(U)}=\inf \{ {\|u\|}_U : u \in U,\ \gamma u=x \},


$$


будет $B$-пространством, изометричным, очевидно,


фактор-пространству $U/(\ker \gamma \cap U)$, причем $\gamma(U) \subset


\Psi$ и $\gamma \in L(U, \gamma(U))$. Кроме того, если


$B$-пространство $X$ непрерывно вложено в $\Psi$, то, для того


чтобы $\gamma$ непрерывно отображало $U$ в $X$, достаточно (и,


разумеется, необходимо), чтобы $\gamma(U)$ было подмножеством


пространства $X$. Если $\gamma$ непрерывно отображает


$B$-пространство $U$ на $B$-пространство $X$, то фактор-норма


пространства $\gamma (U)$ эквивалентна норме пространства~$X$.





Пусть $B$-пространства $U_1$, $U_2$ вложены в ТВП $\Phi$.


Линейные множества $U_1 \cap U_2$ и $U_1+U_2=\{u_1+u_2 : u_j


\in U_j$, $j=1,2 \}$ являются $B$-пространствами относительно норм


(\cite{Trib}, c.\,15)


\begin{align*}


{\|u\|}_{U_1 \cap U_2}&=\|u\|_{U_1}+\|u\|_{U_2}, \\


%


{\|u\|}_{U_1+U_2}&=\inf \{ {\|u_1 \|}_{U_1}+{\|u_2


\|}_{U_2} \}


\end{align*}


($\inf$ берется по всем представлениям $u=u_1+u_2$, $u_j


\in U_j$). Операции (функторы) $\cap$ и $+$


двойственны друг другу в следующем смысле: $(U_1 \cap U_2)^* =


{U_1}^*+{U_2}^*$, $(U_1 +U_2)^*={U_1}^* \cap {U_2}^*$


\cite{Aron}.





Оператор $\gamma \in L(U, X)$ называется ретракцией, если


существует оператор $\beta \in L(X, U)$, называемый


коретракцией такой, что


$$


\gamma \beta x=x \ \ \forall x \in X,


$$


т.\,е. оператор $\gamma$ -- ретракция, если для него существует


правый обратный на $X$, называемый коретракцией. Заметим, что


оператор $\gamma \in L(U, X)$ будет являться ретракцией с


соответствующей коретракцией $\beta \in L(X, U)$ тогда и только


тогда, когда сопряженный к $\beta$ оператор $\beta^* \in L(U^*,


X^*)$ будет ретракцией с соответствующей ему коретракцией


$\gamma^* \in L(X^*, U^*)$. Это следует из того, что произведение


операторов $\gamma\beta \in L(X,X)$ будет тождественным оператором


в~$X$ тогда и только тогда, когда произведение $\beta^*\gamma^* =


(\gamma\beta)^* \in L(X^*,X^*)$ будет тождественным оператором в


сопряженном $X^*$.





Одним из содержательных и важных примеров ретракций в


функциональных пространствах является оператор следа на


некоторое многообразие. Для открытого множества $\Omega \subset


R^n$, $p, q \in (1, \infty)$, $s \in (0,\infty)$, мы используем


стандартные обозначения: через $W^s_p(\Omega)$ и


$B^s_{p, q}(\Omega)$ обозначаются


пространства Соболева--Слободецкого и Бесова


соответственно. Тогда, например, для $\Omega= R^n_+=\{x =


(x_1, x_2, \dotsc, x_n) \in R^n : x_n > 0 \}$, $x'=(x_1, x_2,


\dotsc, x_{n-1})$ отображение $\gamma$, задаваемое формулой


$$


\gamma u=\bigg(u(x',0), \frac{\partial u(x',0)}{\partial x_n}, \dotsc,


\frac{\partial ^m u(x',0)}{\partial x_n^m}\bigg),


$$


является ретракцией пространства $W^s_p (R^n_+)$ на


$\prod \limits_{j=0}^m W_p^{s-1/p-j}(R^{n-1})$, где $m=\max \{k


\in \bold Z : k < s-1/p \}$ (\cite{Trib}, c.\,267).





Для более короткой записи всюду в дальнейшем, когда это не


вызывает недоразумений, для $\gamma \in L(\Phi, \Psi)$ и $U


\subset \Phi$ через $U_0$ обозначается ядро сужения оператора


$\gamma$ на $U$, т.\,е. $U_0 =\ker \gamma \cap U$. Отметим, что


коретракция $\beta$ осуществляет изоморфизм на дополнение


к ядру $U_0$ в пространстве~$U$ (\cite{Trib}, c.\,21) и пространство $U$


представляется алгебраической и топологической прямой суммой $U


= U_0 \oplus \beta(X)$. Дополняемость ядра есть характеризущее


свойство ретракции, как показывает следующая





\begin{lems} \label{one}


Оператор $\gamma \in L(U, X)$


является ретракцией тогда и только тогда,


когда $U_0$ дополняемо в $U$ и область значений $\gamma(U)=X$.


Если это выполнено и $\pi$ --- оператор проектирования на $U_0$,


то отображение $\Pi$, определяемое формулой


\begin{equation}


\Pi u=(\pi u, \gamma u), \label{Pi}


\end{equation}


осуществляет изоморфизм пространства $U$ на $U_0


\times X$.


\end{lems}





\begin{pf}


{\bf Необходимость.} Из определения


ретракции следует, что $\gamma(U)=X$. Пусть $\beta : X


\to U$ --- коретракция для $\gamma$. Определим оператор


$\pi \in L(U, U)$ формулой $ \pi u=u - \beta \gamma u$. Тогда


$\gamma \pi u=\gamma u - \gamma \beta \gamma


u=\gamma u - \gamma u=0$, т.\,е. $\pi u \in U_0$; кроме того, для


$u \in U_0$ \ $ \pi u=u - \beta \gamma u=u$, откуда следует, что $\pi$


является проектором на $U_0$.


\smallskip





{\bf Достаточность.} Покажем сначала, что отображение $\Pi$,


определенное формулой (\ref{Pi}), осуществляет изоморфизм


пространства $U$ на $U_0 \times X$. Очевидно, оператор $\Pi$


непрерывен. Если $Tu=0$, то $\gamma u=0$, т.\,е. $u \in U_0$; с


другой стороны, $0=\pi u=u$, откуда следует, что $\Pi$


взаимнооднозначен. Пусть теперь $(v, x) \in U_0 \times X$ ---


произвольный элемент. Так как $\gamma(U)=X$, то найдется элемент


$w \in U$ такой, что $\gamma w=x$. Положим $ u=v+w - \pi w$.


Тогда $\pi u=\pi v+\pi w - \pi^2 w=\pi v=v$ и $ \gamma


u=\gamma v+\gamma w - \gamma \pi w=\gamma w=x$, т.\,е. $(v,


x)=\Pi u$. Тем самым установлено, что отображение $\Pi$ есть


отображение ``на''. Итак, оператор $\Pi$ взаимнооднозначно и


непрерывно отображает $U$ на $U_0 \times X$. По теореме Банаха об


открытом отображении обратный оператор $\Pi^{-1}$ также


непрерывен, т.\,е. $\Pi$ --- изоморфизм.





Положим теперь $\beta x=\Pi^{-1} (0,x)$ для каждого $x \in X$.


Очевидно, определяемый таким образом оператор


$\beta$ линеен, непрерывен и, кроме того,


$(0, x)=\Pi \beta x=(\pi \beta x, \gamma \beta x)$,


т.\,е. $\gamma \beta x=x$, откуда по определению следует,


что оператор $\gamma$ является ретракцией.


\end{pf}





\begin{note}


Из леммы, в частности, вытекает, что


топологические и структурные свойства $B$-пространства $U$


такие, как сепарабельность, рефлексивность, существование базиса


Шаудера, изоморфность гильбертову пространству и т.\,п.,


определяются соответствующими свойствами ядра ретракции и ее


области значений.


\end{note}





\begin{note}


Если полунорма $p$ непрерывна на $U$ и эквивалентна норме


пространства $U$ на ядре ретракции $U_0$, то норма


$$


u \rightarrow p(u)+{\|\gamma u \|}_X


$$


эквивалентна норме пространства $U$ на всем $U$. Это


следует из свойств отображения (\ref{Pi}).


\end{note}





\begin{note}


Если $\gamma \in L(U,X)$ и $U_0$ дополняемо в $U$,


то оператор $\gamma$ является ретракцией $U$ на $\gamma(U)$. В


частности, если оператор $\gamma$ имеет конечномерное ядро, т.е.


$\dim U_0 < \infty$, то $\gamma$ является ретракцией $U$ на


$\gamma(U)$.


\end{note}





\begin{note}


Если пространство $U$ изоморфно гильбертову пространству,


то любое замкнутое подпространство этого пространства дополняемо в


$U$. Из доказанного утверждения следует, что в этом случае любое


непрерывное отображение из $U$ на некоторое $B$-прост\-ран\-ст\-во~$X$


будет являться ретракцией, причем само это пространство $X$ с


необходимостью должно быть изоморфно гильбертову пространству.


\end{note}





Всюду далее в этом пункте будем предполагать, что $\gamma$


есть ретракция $U$ на $X$. Для $u \in U$, $u^* \in U^*$ полагаем


$\la u,u^*\ra=\la u,u^* \ra_U=u^*(u)$.





\begin{lems}


Для любого линейного непрерывного функционала $u^* \in


U^*$ существует единственная пара $v^* \in {U_0}^*$, $x^* \in


X^*$ такая, что


\begin{equation}


\la u, u^*\ra _U=\la \pi u, v^*\ra _{U_0}+\la \gamma u,


x^*\ra _X \ \ \forall u \in U. \label{func}


\end{equation}


Очевидно, справедливо и обратное\/${:}$ каждая пара $(v^*,


x^*) \in {U_0}^* \times X^*$ формулой $(\ref{func})$


определяет линейный непрерывный функционал $u^* \in U^*$.


\end{lems}





\begin{pf}


Сопряженное к декартовому произведению


$U_0 \times X$ естественным образом отождествим с ${U_0}^*


\times X^*$ отношением двойственности


$$


\la (v,x), (v^*, x^*)\ra _{U_0 \times X}=\la v, v^*\ra _{U_0}+\la x,


x^*\ra _X


$$


для $v \in U_0$, $x \in X$, $v^* \in {U_0}^*$, $x^*


\in X^*$. По лемме \ref{one} оператор в формуле (\ref{Pi}) является


изоморфизмом $U$ на $U_0 \times X$, следовательно, сопряженный


оператор $\Pi^*$ является изоморфизмом ${U_0}^* \times X^*$ на


$U^*$. С каждым элементом $u^* \in U^*$ формулой


$(v^*, x^*)={\Pi^*}^{-1}u^*$


свяжем пару $(v^*, x^*) \in {U_0}^* \times X^*$.


Тогда для $u \in U$ имеем следующую цепочку равенств:


\begin{multline*}


\la u, u^*\ra _U=\la u, \Pi^*(v^*, x^*)\ra _U=\la Tu, (v^*,


x^*)\ra _{U_0 \times X} = \\


%


= \la (\pi u, \gamma u), (v^*, x^*)\ra _{U_0 \times X}


= \la \pi u, v^*\ra _{U_0}+\la \gamma u, x^*\ra _X,


\end{multline*}


что доказывает как разложение (\ref{func}), так и


его единственность.


\end{pf}





\begin{notenonum}


Из формулы (\ref{func}) следует, что для $u \in U_0$


$$


\la u, u^*\ra _U=\la \pi u, v^*\ra _{U_0}=\la u, v^*\ra _{U_0},


$$


т.\,е. функционал $v^* \in {U_0}^*$ есть сужение


функционала $u^* \in U^*$ на подпространство $U_0$.


\end{notenonum}





В следующей теореме дается критерий плотности подмножества


в терминах ядра ретракции и ее области значений.





\begin{thms} \label{critdens}


Для того чтобы линейное множество $L \subset U$


было плотно в $U$, необходимо и достаточно, чтобы замыкание $L$


содержало ядро $U_0$ и множество $\gamma(L)$ было плотно в $X$.


\end{thms}





\begin{pf} {\bf Необходимость.} Так как


замыкание $L$ совпадает с $U$, то оно содержит в качестве


подпространства ядро $U_0$. Далее, по лемме \ref{one} оператор $\Pi$,


определяемый формулой (\ref{Pi}), есть изoморфизм $U$ на $U_0


\times X$, следовательно, множество $\Pi(L)$ плотно в последнем


пространстве, следовательно, $\gamma(L)$ плотно в $X$.


\smallskip





{\bf Достаточность.} Пусть $M$ --- замыкание $L$ в $U$.


Нужно показать, что $M=U$; последнее равносильно тому (в силу


замкнутости $M$), что если $u^* \in U^*$ --- любой функционал такой,


что $M \subset \ker u^*$, то $u^*=0$. Итак, пусть $u^* \in U^*$


и $M \subset \ker u^*$. По предыдущей лемме для $u^*$ найдется пара


элементов $(v^*, x^*) \in {U_0}^* \times X^*$, для которой будет


справедлива формула (\ref{func}). Так как по условию $U_0


\subset M$, то, подставляя в эту формулу элементы $u \in U_0$,


получим


$\la u, v^*\ra _{U_0}=0$\; $\forall u \in U_0$,


т.\,е. $v^*=0$. Но тогда


$\la \gamma u, x^*\ra _X=0$\; $\forall u \in L \subset M$.


В силу плотности $\gamma(L)$ в $X$ из последнего


тождества заключаем, что $x^*=0$. Из представления (\ref{func})


теперь следует $u^*=0$.


\end{pf}





\begin{lems} \label{denseps}


Пусть линейное множество $L \subset U$ плотно в $U$ и $L \bigcap U_0$


плотно в $U_0$. Если $u \in U$ и $\gamma u \in \gamma(L)$, то


для произвольного $\varepsilon > 0$ найдется $u_\varepsilon \in L$ такой, что


$$


{\left\| {u-u_\varepsilon } \right\|}_{U}


\le \varepsilon \ \ \text{и}\ \ \gamma u_\varepsilon


=\gamma u.


$$


\end{lems}





{\bf Доказательство.} 


Пусть элемент $u \in U$ и $\gamma u \in \gamma(L)$, т.\,е.


найдется $w \in L$, что $\gamma w=\gamma u$.


Тогда $v=u-w \in U_0$. Для произвольного $\varepsilon > 0$ найдем


элемент $v_\varepsilon \in L \bigcap U_0$ такой, что


${\left\| {v-v_\varepsilon } \right\|}_U \le


\varepsilon$. Положим $u_\varepsilon=v_\varepsilon+w \in L$. По построению


$$


{\left\| {u-u_\varepsilon} \right\|}_U=


{\left\| {v-v_\varepsilon } \right\|}_U \le \varepsilon \quad


\text{и} \quad \gamma u_\varepsilon=\gamma v_\varepsilon +


\gamma w=\gamma w=\gamma u.\quad\square


$$


\medskip





В качестве иллюстрации к лемме рассмотрим пространство Соболева


$W_p^1(\Omega)$, $1<p<\infty$, и $\Gamma=\partial \Omega \in


C^\infty$ ($\Omega$ --- ограниченное множество). Пусть $\gamma u =


u|_{\Gamma}$. Оператор $\gamma$ является ретракцией из


$W_p^1(\Omega)$ на $W_p^{1-1/p}(\Gamma)$. Как известно, класс


функций $C^{\infty}(\overline{\Omega})$ плотен в $W_p^1(\Omega)$,


а множество финитных функций $C_0^{\infty}(\Omega) \subset


C^{\infty}(\overline{\Omega}) \bigcap \ker \gamma$ плотно в


подпространстве $\overset{\circ}{W}{}_p^1(\Omega)=W_p^1(\Omega)


\bigcap \ker \gamma$. Пусть функция $u \in W_p^1(\Omega)$ и


$\gamma u \in C^\infty(\Gamma)$. По теореме для $u|_\Gamma=\gamma


u$ найдется сколь угодно близкое к $u$ в метрике пространства


$W_p^1(\Omega)$ продолжение с $\Gamma$ на $\Omega$ из класса


$C^\infty(\overline{\Omega})$.








\section{$B$-пространства с нормой графика} \label{force}





Как и в п.\,1, здесь также предполагается, что оператор


$\gamma \in L(\Phi, \Psi)$. Пусть $U \subset \Phi$ и $X \subset


\Psi$ --- два $B$-пространства. Определим пространство


$$


\gamma (U, X) \overset{\mathrm{def}}{=}{} \{u \in U: \gamma u \in X\},


$$


наделенное нормой графика


\begin{equation}


{\|u\|}_{\gamma(U,X)}={\|u\|}_U+{\|\gamma u\|}_X.


\label{normaV}


\end{equation}


Из определения нормы (\ref{normaV}) и полноты


пространств $U$, $X$ с очевидностью вытекает, что $\gamma(U,X)$


является $B$-пространством, непрерывно вложенным в пространство


$U$, причем $\gamma (U,X)=U$ тогда и только тогда, когда


$\gamma(U)$ является подмножеством пространства $X$.





Далее установим


некоторые общие свойства пространств $\gamma(U,X)$ в терминах


"усиливаемого" пространства $U$ и ``усиливающего'' $X$ при


фиксированном отображении $\gamma \in L(\Phi, \Psi)$.





\begin{thms} \label{t1}


Для пространства $V=\gamma(U, X)$


справедливы следующие утверждения\/${:}$


\begin{itemize}


\item[(i)] оператор $\gamma$ непрерывно отображает $V$ на


$X \cap \gamma(U);$


\item[(ii)] отображение $u \rightarrow (u, \gamma u)$ осуществляет


изометрию $V$ на замкнутое подпространство пространства $U \times


X$.


\end{itemize}


Если, кроме того, $U_0=\ker \gamma \cap U$ дополняемо


в $V$, то


\begin{itemize}


\item[(iii)] сужение $\gamma$ на $V$ есть ретракция $V$ на


$X \cap \gamma(U);$


\item[(iv)] отображение $(\ref{Pi})$, где $\pi$ --- проектор на


$U_0$, есть изоморфизм $V$ на $U_0 \times (X \cap \gamma(U));$


\item[(v)] $V$ плотно в $U$ тогда и только тогда, когда $X \cap


\gamma(U)$ плотно в $\gamma(U)$.


\end{itemize}


\end{thms}





\begin{pf}


Утверждения (i), (ii) непосредственно


вытекают из определения нормы (\ref{normaV}).





(iii), (iv). В силу (i) \ $\gamma \in L(V, X \cap \gamma(U))$,


причем $\gamma(V)=X \cap \gamma(U)$. Из леммы \ref{one} следует


теперь, что $\gamma$ есть ретракция $V$ на $X \cap \gamma(U)$ и


отображение (\ref{Pi}) есть изморфизм $V$ на $U_0 \times (X \cap


\gamma(U))$.





Утверждение (v) следует из теоремы \ref{critdens}, т.\,к.


$V \supset U_0$ и $\gamma(V)=X \cap \gamma(U)$.


\end{pf}





\begin{notenonum}


Если подпространство $U_0 \subset V$


дополняемо в $U$, то оно дополняемо и в пространстве $V$.


\end{notenonum}





\begin{cor}


Пусть $V \subset U$, $X \subset \Psi$.


Для того чтобы $V=\gamma (U, X)$, необходимо и достаточно,


чтобы сужение $\gamma$ на $V$ непрерывно отображало $V$ на


$X \cap \gamma(U)$ и было выполнено включение $U_0 \subset V$.


\end{cor}





\begin{pf} {\bf Необходимость} следует из теоремы.


\smallskip





{\bf Достаточность.} Из условия следует, что $V \subset


\gamma (U,X)$. Пусть $u \in \gamma (U,X)$. Так как $\gamma(V) =


X \cap \gamma (U)$, то найдется такой элемент $v \in V$, что


$\gamma u=\gamma v$. Но тогда $u-v \in \ker \gamma \cap U=U_0


\subset V$. Отсюда $u=v+(u-v) \in V$. В силу


произвольности $u$ заключаем, что множества $V$ и $\gamma


(U,X)$ совпадают между собой. Так как норма усиленного


пространства $\gamma(U,X)$ слабее нормы $V$, то из теоремы


Банаха об открытом отображении следует, что нормы пространств


$V$ и $\gamma(U,X)$ эквивалентны.


\end{pf}





\begin{cor}


Пусть $V=\gamma (U,X)$. Для того


чтобы подмножество $K \subset V$ было относительно компактным в


$V$, необходимо и достаточно, чтобы $K$ было относительно


компактным в $U$ и $\gamma (K)$ было относительно компактным в $X$. 


\end{cor}





{\bf Доказательство} следует из утверждения (ii) теоремы \ref{t1}.


\medskip





Из утверждения (ii) теоремы \ref{t1} вытекает также, что если


пространства $U$ и $X$ оба обладают одним из таких свойств, как


сепарабельность, рефлексивность, изоморфность гильбертову


пространству, то соответствующим свойством будет обладать и


пространство $V=\gamma (U,X)$.$\quad\square$





\begin{cor}


Пусть $U_0$ дополняемо в пространстве


$V=\gamma(U,X)$. Для того чтобы линейное множество $L$


было плотно в пространстве $V$, необходимо и достаточно, чтобы


замыкание $L$ в $V$ содержало подпространство $U_0$ и


множество $\gamma(L)$ было плотно в прoстранстве $X \bigcap \gamma(U)$.


\end{cor}





\begin{pf}


По теореме оператор $\gamma$ является


ретракцией пространства $V$ на пространство $Y=X \bigcap \gamma(U)$.


Поэтому утверждение немедленно следует из теоремы \ref{critdens}.


\end{pf}





\begin{examnonum}


Пусть $\Omega \subset R^n$ --- открытое


множество с границей $\partial\Omega \in C^{\infty}$. Тогда


оператор следа $\gamma u \equiv u|_{\partial\Omega}$ является


ретракцией пространства $W^1_p(\Omega)$ на


$W_p^{1-1/p}(\partial\Omega)$ (\cite{Necas}, c.\,103). Пусть $q \in (1,


\infty)$, $q \geq (n-1)p/n$. Тогда


$W_q^1(\partial\Omega)$ непрерывно и плотно вложено в


$W_p^{1-1/p}(\partial\Omega)$. Из доказанной


выше теоремы следует, что усиленное пространство Соболева


$$


V_{p, q}^{1, 1} \equiv \gamma(W_p^1(\Omega),


W_q^1(\partial\Omega)) =


\{u \in W_p^1(\Omega): \gamma u \in W_q^1(\partial\Omega) \}


$$


будет сепарабельным рефлексивным $B$-пространством,


непрерывно и плотно вложенным в пространство Соболева


$W_p^1(\Omega)$, причем оператор следа на $\partial\Omega$


является ретракцией $V_{p, q}^{1, 1}$ на


$W_q^1(\partial\Omega)$. Если $p=q=2$, то усиленное пространство


Соболева $V_{2, 2}^{1, 1}$ будет гильбертовым пространством.


\end{examnonum}





\begin{thms}


Пусть $\gamma \in L(U,X)$ --- ретракция и $\beta \in L(X,


U)$ --- соответствующая $\gamma$ коретракция. Если $B$-пространство


$Y \subset \Psi$, то для пространства $V=\gamma(U,Y)$ имеет


место двойственная формула


$$


U^*=\beta^*(V^*, X^*),


$$


т.\,е. $U^*$ есть усиление $V^*$ полунормой


${\|\beta^*v^*\|}_{X^*}$.


\end{thms}





{\bf Доказательство.} Как было отмечено в п.\,1,


сопряженный к коретракции $\beta$ оператор


$$


\beta^* \in L(U^*, X^*) \cap L(V^*, Y^*)


$$


является ретракцией как из $U^*$ на $X^*$, так и из


$V^*$ на $(X \cap Y)^*=X^*+Y^*$ (в силу утверждения (ii) теоремы


\ref{t1}). Поэтому для $u^* \in U^*$ имеем


$$


{\|u^*\|}_{\beta ^*(V^*,X^*)} =


{\|u^*\|}_{V^*}+{\|\beta^*u^*\|}_{X^*} \leq c{\|u^*\|}_{U^*}.


$$


Отсюда следует включение $U^* \subset \beta^*(V^*,


X^*)$. Пусть $u^* \in \beta^*(V^*, X^*)$ --- произвольный элемент,


т.\,е. $\beta^* u^* \in X^*$. Нужно установить, что $u^* \in U^*$.


Формула


$\pi u=u - \beta\gamma u$


определяет проектор на $U_0$. Через $v^* \in {U_0}^*$


обозначим сужение $u^*$ на $U_0$. Тогда для $u \in U$ имеем


$$


\la u, u^*\ra _V=\la \pi u, u^*\ra _V+\la \beta\gamma u, u^*\ra _V=


\la \pi u, v^*\ra _{U_0}+\la \gamma u, \beta^* u^*\ra _{X \cap Y}.


$$


Отсюда (т.\,к. $\beta^* u^* \in X^*$) вытекает


непрерывность $u^*$ в топологии пространства $U$, т.\,е. $u^* \in


U^*$. В силу произвольности $u^* \in \beta^*(V^*, X^*)$ это


доказывает теорему.





\begin{thms} \label{tdens}


Пусть $W$ непрерывно и плотно вложено в $U$,


$\gamma$ является ретракцией на $W$ и на $U$ и пусть $X=\gamma(W)$. Тогда


$W$ плотно в $V=\gamma(U,X)$ и $V$ плотно в $U$, причем для любого $u \in V$


и произвольного $\varepsilon > 0$ найдется $u_\varepsilon \in W$ такой, что


$$


{\left\| {u-u_\varepsilon } \right\|}_{V}


\le \varepsilon \ \ \text{и} \ \ \gamma u_\varepsilon=\gamma u.


$$


\end{thms}





\begin{pf}


Так как $W \subset U$, то $X \subset \gamma(U)$,


а потому $X=\gamma(V)$ (утверждение (i) теоремы \ref{t1}).


Пусть $\beta \in L(X,W)$ --- коретракция для $\gamma \in L(W,X)$, т.\,е.


$$


\gamma \beta x=x \ \ \forall x \in X.


$$


Но тогда отсюда и из равенства $X=\gamma(V)$ по определению


следует, что $\gamma$ есть ретракция из $V$ на $X$ с соответствующей


коретракцией $\beta \in L(X,V)$. 


Поскольку $W$ плотно в $U$, то $X=\gamma(W)$ плотно


в $\gamma(U)$. По утверждению (v) теоремы \ref{t1}


пространство $V$ плотно в $U$.





Положим для $u \in U$ \ $\pi u=u - \beta \gamma u$.


Оператор $\pi$ является проектором в $W$ на $W \bigcap \ker \gamma$,


а также в $V$ и $U$ на $V \bigcap \ker \gamma=U \bigcap \ker \gamma$.


Следовательно, $\pi(W)=W \bigcap \ker \gamma$


плотно в $\pi(V)=\pi(U)=V \bigcap \ker \gamma$. Возьмем


произвольный элемент $u \in V$.


Найдется такое $w \in W$, что $\gamma w=\gamma u$.


Тогда $v=u-w \in V \bigcap \ker \gamma$.


Для произвольного $\varepsilon > 0$ найдем


$v_\varepsilon \in W \bigcap \ker \gamma$


такой, что ${\left\| {v-v_\varepsilon } \right\|}_V \le


\varepsilon$. Положим $u_\varepsilon=v_\varepsilon+w \in W$. По построению


${\left\| {u-u_\varepsilon} \right\|}_V=


{\left\| {v-v_\varepsilon } \right\|}_V \le \varepsilon$


и


$\gamma u_\varepsilon=


\gamma v_\varepsilon+\gamma w=\gamma w=\gamma u$.


\end{pf}








\begin{cornonum}


Пусть $W$ непрерывно и плотно вложено в $U$ и задано конечное


семейство линейных непрерывных 


функционалов $u_i^* \in U^*$, $i=\overline{1,m}$.


Тогда для любых $u \in U$ и $\varepsilon > 0$


найдется $u_\varepsilon \in W$ такой, что


$$


{\left\| {u-u_\varepsilon } \right\|}_{U}


\le \varepsilon \ \ \text{и}\ \ {\la u_\varepsilon, u_i^*\ra }_U


= {\la u, u_i^*\ra }_U \ \ \forall i=\overline{1,m} .


$$


\end{cornonum}





\begin{pf}


Не ограничивая общности, будем считать, что функционалы $u_i^*$


образуют линейно независимое множество. Формулой


$$


\gamma u={({\la u, u_i^*\ra }_U)}_{i=1,m}


$$


определим линейный непрерывный оператор из $U$ на $R^m$. 


В силу конечномерности,


$\gamma$ является ретракцией на $W$ и на $U$. По теореме для любых


$u \in V=\gamma(U,R^m)=U$ и $\varepsilon >0$


найдется $u_\varepsilon \in W$ такой, что


${\left\| {u-u_\varepsilon } \right\|}_{U} \le \varepsilon$


и


$\gamma u_\varepsilon= \gamma u$.


\end{pf}








В следующей теореме устанавливается интерполяционное


свойство усиленных пространств.





\begin{thms} \label{tinterp}


Пусть $U_j \subset


\Phi$, $X_j \subset \gamma(U_j)$ и $U_j \cap \ker \gamma$


дополняемо в $U_j$ для $j=1,2$. Тогда для произвольного


интерполяционного функтора $F$ имеет место равенство


$$


F\{\gamma(U_1,X_1), \gamma(U_2,X_2) \} =


\gamma (F\{U_1,U_2\}, F\{X_1,X_2\}).


$$


\end{thms}





Доказательству теоремы предпошлем два вспомогательных утверждения.





\begin{lems} \label{proector}


Пусть оператор $\pi \in L(U,U)$ --- проектор


$($т.\,е. $\pi^2=\pi)$. Если $V \subset U$ и $\pi(V) \subset V$, то


$\pi \in L(V,V)$ является оператором проектирования на


подпространство $\pi(V)=V \cap \pi(U)$.


\end{lems}





\begin{pf}


Ясно, что $\pi(V) \subset V \cap


\pi(U)$. С другой стороны, если $v \in V \cap \pi(U)$, то


найдется элемент $u \in U$ такой, что $v=\pi u$. Тогда


$v=\pi u=\pi^2 u=\pi(\pi u) \in \pi(V)$.


Это доказывает, что $\pi(V)=V \cap \pi(U)$.


\end{pf}





\begin{lems} \label{lemmainterp}


Пусть $U_j \subset \Phi$ и $\pi \in


L(U_j,U_j)$ является проектором в $U_j$ $(j=1,2)$. Положим


$W=\pi(U_1+U_2)$. Тогда для произвольного интерполяционного


функтора $F$ имеет место равенство


$F\{U_1 \cap W, U_2 \cap W \}=F\{U_1,U_2\} \cap W$.


\end{lems}





Доказательство утверждения можно найти, например, в


(\cite{Trib}, c.\,138).





\begin{pf*}{Доказательство теоремы \ref{tinterp}}


Введем обозначения


\begin{gather*}


U=F\{U_1,U_2\},\ \ X=F\{X_1,X_2\},\ \ Y_j=\gamma(U_j),\ \


Y=F\{Y_1,Y_2\}, \\


V_j=\gamma(U_j,X_j),\ \ j=1,2,\ \ V=F\{V_1,V_2\}.


\end{gather*}


Необxодимо показать, что $V=\gamma(U,X)$.


Поскольку $V_j \subset U_j$, то $V \subset U$. Далее,


т.\,к. $\gamma \in L(V_j, X_j)$ --- ретракция (утверждение (iii)


теоремы \ref{t1}), то $\gamma$ является также ретракцией, как $U$


на $Y$, так и $V$ на $X \subset Y$ (\cite{Trib}, c.\,21). По следствию


1 теоремы \ref{t1} осталось показать, что $\ker \gamma \cap U


\subset V$.


Положим $\pi u=u - \beta\gamma u$ для $u \in U_1+U_2$,


где $\beta \in L(Y_1+Y_2, U_1+U_2)$ --- коретракция для


$\gamma$. Как видно из доказательства леммы \ref{one}, оператор


$\pi$ проектирует $U_1+U_2$ на $W=(U_1+U_2) \cap \ker \gamma$, а


также каждое $U_j$ на $U_j \cap \ker \gamma$ и $U$ на $U \cap \ker


\gamma$. Но согласно леммe \ref{proector} \ $U_j \cap \ker \gamma=U_j


\cap W$ и $U \cap \ker \gamma=U \cap W$. Теперь по лемме


\ref{lemmainterp} получаем


$$


U \cap \ker \gamma=U \cap W=F\{U_1 \cap W,U_2 \cap W\}.


$$


Поскольку $U_j \cap W=U_j \cap \ker \gamma \subset


V_j$, то $U \cap \ker \gamma=F\{U_1 \cap W, U_2 \cap W \}


\subset F\{V_1, V_2 \}=V$.


\end{pf*}





Установим условия компактного вложения одного усиленного


пространства в другое.





\begin{lems} \label{lcomp}


Если пространство $U_2$ компактно


вложено в $U_1$ $($т.\,е. единичный шар пространства $U_2$


относительно компактен в пространстве $U_1)$, то пространство


$\gamma(U_2)$ компактно вложено в $\gamma(U_1)$.


\end{lems}





{\bf Доказательство.} Пусть $K=\{u \in U_2:{\|u\|}_{U_2} < 1\}$.


Множество $K$ открыто в $U_2$, с другой стороны, $K$


относительно компактно в $U_1$. По теореме Банаха об открытом


отображении (см., напр., \cite{Iosida}, c.\,112) оператор $\gamma$


является открытым отображением $U_2$ на $\gamma(U_2)$. Поэтому


$\gamma(K)$ является окрестностью нуля в $\gamma(U_2)$. С другой


стороны, непрерывный оператор переводит относительно компактные


множества в относительно компактные. Поэтому $\gamma(K)$


относительно компактно в $\gamma(U_1)$. Таким образом,


$\gamma(K)$ является окрестностью нуля в пространстве


$\gamma(U_2)$, которая относительно компактна в пространстве


$\gamma(U_1)$. Это доказывает компактность тождественного


отображения из $\gamma(U_2)$ в $\gamma(U_1)$, т.\,е. компактность


вложения $\gamma(U_2)$ в $\gamma(U_1)$.





\begin{thms} \label{tcomp}


Пусть $U_2 \subset U_1 \subset \Phi$,


$X_1, X_2 \subset \Psi$. Введем обозначения\/${:}$ $V_j =


\gamma(U_j,X_j)$, $U_{j,0}=U_j \cap \ker \gamma$. Справедливы


следующие утверждения\/${:}$


\begin{itemize}


\item[(i)] если $U_2$ компактно вложено в $U_1$ и $X_2 \cap


\gamma(U_2)$ компактно вложено в $X_1$, то $V_2$ компактно


вложено в $V_1;$


\item[(ii)] если $V_2$ компактно вложено в $V_1$, то $U_{2,0}$


компактно вложено в $U_{1,0}$ и $X_2 \cap \gamma(U_2)$ компактно


вложено в $X_1 \cap \gamma(U_1);$


\item[(iii)] если каждое из $U_{j,0}$ дополняемо в $V_j$, то


утверждение {\em (ii)} обратимо.


\end{itemize}


\end{thms}





\begin{pf}


Положим $Y_j=X_j \cap \gamma(U_j)$.





(i). Пусть $K$ --- единичный шар пространства $V_2$. Тогда


$K$ ограничено в $U_2$ и, следовательно, относительно компактно


в $U_1$. Далее, $\gamma(K)$ ограничено в $Y_2$ и потому, как


следует из условия, относительно компактно в $X_1$. Из


следствия 2 теоремы \ref{t1} теперь вытекает относительная


компактность $K$ в $V_1$.





(ii). Так как $U_{j,0}$ замкнуто в $V_j$, то из компактности


вложения $V_2$ в $V_1$ следует компактность вложения $U_{2,0}$ в


$U_{1,0}$. Поскольку $\gamma(V_j)=Y_j$, то из леммы \ref{lcomp}


следует, что $Y_2$ компактно вложено в $Y_1$.





Утверждение (iii) является очевидным следствием утверждения


(iv) теоремы \ref{t1}.


\end{pf}





\begin{cornonum}


Пространство $\gamma(U,X)$ компактно


вложено в $U$ тогда и только тогда, когда подпространство $U_0=U


\cap \ker \gamma$ конечномерно и $X \cap \gamma(U)$ компактно


вложено в $\gamma(U)$.


\end{cornonum}





\begin{pf}


В силу утверждения (ii) теоремы, из


компактности вложения $\gamma(U,X)$ в $U=\gamma(U, \gamma(U))$


вытекает, что, во-первых, $U_0$ компактно вкладывается само в


себя, что может быть только в случае его конечномерности;


во-вторых, $X \cap \gamma(U)$ компактно вкладывается в


$\gamma(U)$.


\end{pf}





\begin{thms} \label{DeniLions}


Рассмотрим следующие утверждения\/${:}$


\begin{itemize}


\item[1)] $\gamma(U,X)$ компактно вложено в $U;$


\item[2)] $X \cap \gamma(U)$ компактно вложено в $\gamma(U);$


\item[3)] множество $X \cap \gamma(U)$ замкнуто в $X;$


\item[4)] на множестве $X \cap \gamma(U)$ норма пространства


$\gamma(U)$ слабее нормы $X$, т.\,е. для некоторой постоянной


$c>0$ имеет место неравенство


$$


{\|x\|}_{\gamma(U)} \leq c {\|x\|}_X \ \ \forall x \in X


\cap \gamma(U);


$$


\item[5)] для некоторой постоянной $c >0$ справедливо неравенство


$$


\inf \{ {\|u+v\|}_U : v \in U_0 \} \leq c {\| \gamma u\|}_X


\ \ \forall u \in \gamma(U,X);


$$


\item[6)] для некоторой постоянной $c >0$ справедливо неравенство


$$


\inf \{ {\|u+v\|}_{\gamma(U,X)} : v \in U_0 \} \leq c {\|


\gamma u\|}_X \ \ \forall u \in \gamma(U,X).


$$


\end{itemize}


Тогда имеют место импликации $1) \Rightarrow 2) \Rightarrow 3)


\Leftrightarrow 4) \Leftrightarrow 5) \Leftrightarrow 6)$.


\end{thms}








\begin{lems} \label{XYcompY}


Если $B$-пространства $X, Y \subset


\Psi$ таковы, что $X \cap Y$ компактно вложено в $Y$, то $X \cap


Y$ замкнуто в $X$, или, что равносильно, на множестве $X \cap Y$


норма $Y$ слабее нормы~$X$.


\end{lems}





{\bf Доказательство.} Нужно показать, что


$$


{\|x\|}_Y \leq c {\|x\|}_X \ \ \forall x \in X \cap Y.


$$


Если это не так, то для каждого натурального $n \ge 1$


найдется элемент $x_n \in X \cap Y$ такой, что


${\|x_n\|}_Y > n {\|x_n\|}_X$.


Положим $y_n=x_n / {\|x_n\|}_Y$. Тогда для каждого $n


\ge 1$ \ ${\|y_n\|}_Y=1$ и ${\|y_n\|}_X < 1/n$, откуда


следует, что последовательность $(y_n)$ ограничена в $X \cap Y$.


Разрeжая, если это необходимо, последовательность $(y_n)$,


добьемся ее сходимости в $Y$ к некоторому элементу $y \in Y$


(это можно сделать в силу компактности вложения $X \cap Y$ в $Y$).


Так как ${\|y_n\|}_Y=1$, то и ${\|y\|}_Y=1$. С другой стороны,


$y_n \to 0$ в $X$, а потому $y_n \to 0$ в $\Psi$. Так как $Y$


непрерывно вложено в $\Psi$, то $y=0$, что противоречит


равенству ${\|y\|}_Y=1$. Полученное противоречие доказывает


лемму.








\begin{pf*}{Доказательство теоремы \ref{DeniLions}}


Импликация $1)


\Rightarrow 2)$ установлена в следствии теоремы \ref{tcomp}.


Импликация $2) \Rightarrow 3)$ получается из леммы \ref{XYcompY}.





Докажем, что 3) равносильно 4). Замкнутость $B$-пространства


$X \cap \gamma(U)$ в $X$ равносильна тому, что на $X \cap


\gamma(U)$ нормы ${\|x\|}_{\gamma(U)}+{\|x\|}_X$ и ${\|x\|}_X$


эквивалентны, что в свою очередь равносильно~4).





Эквивалентность $4) \Leftrightarrow 5)$ вытекает из равенств


$\inf \{ {\|u+v\|}_U : v \in U_0 \}={\| \gamma u


\|}_{\gamma(U)}$\; $\forall u \in U$


и $\gamma(\gamma(U,X))=X \cap \gamma(U)$.





Эквивалентность $5) \Leftrightarrow 6)$ следует из


определения нормы пространства $\gamma(U,X)$.


\end{pf*}





Как показывает следующий пример, теорема \ref{DeniLions}


обобщает на абстрактные $B$-пространства хорошо известную в


теории пространств Соболева теорему Дени--Лионса.





\begin{examnonum}


Пусть $\Omega \subset R^n$ --- область со


свойством конуса. Для мультииндекса $i=(i_1, i_2, \dotsc, i_n)$


через $D^i$ обозначим оператор дифференцирования порядка


$|i|=i_1+i_2+\dotsb+i_n$. Линейное отображение $\gamma u =


(D^i u)_{|i|=m}$ ($m$ --- натуральное число) непрерывно как


отображение из пространства $D'(\Omega)$ в $D'(\Omega)^N$, где


$N$ --- мощность множества мультииндексов $i$ таких, что $|i|=m$.


Пространство Соболева $W^m_p(\Omega)$ ($1<p<\infty$) можно


рассматривать как усиленное пространство $\gamma(U,X)$, где


$U=L_p(\Omega) \subset \Phi=D'(\Omega)$ и $X=L_p(\Omega)^N


\subset \Psi=D'(\Omega)^N$. Заметим, что ядро $\ker \gamma$


совпадает с пространством полиномов $P_m(\Omega)$ степени меньшей $m$


по совокупности переменных (\cite{Sobolev}, c.\,64), следовательно,


конечномерно и дополняемо в $L_p(\Omega)$. Из компактности


вложения $W^m_p(\Omega)=\gamma(U,X)$ в $L_p(\Omega)=U$


следует (в силу импликации $1) \Rightarrow 6)$ теоремы


\ref{DeniLions}) оценка


$$


\inf \{ {\|u+v\|}_{W^m_p(\Omega)} : v \in P_m(\Omega) \}


\leq c \sum_{|i|=m} {\|D^i u\|}_{L_p(\Omega)} \ \ \forall u


\in W^m_p(\Omega).


$$


Это неравенство и является содержанием теоремы Дени--Лионса.


\end{examnonum}





\begin{cornonum}


Пусть $Z$ --- некоторое


нормированное пространство и $A:\gamma(U,X) \to Z$ --- непрерывный


оператор, причем $Au=0$ при $u \in U_0$. Если $X \cap


\gamma(U)$ замкнуто в $X$, то справедлива оценка


$$


{\|Au\|}_Z \le c {\|\gamma u\|}_X \ \ \forall u \in


\gamma(U,X).


$$


\end{cornonum}





\begin{pf}


Для произвольного элемента $v \in U_0$


имеем при фиксированном $u \in \gamma(U,X)$


$$


{\|Au\|}_Z={\|A(u+v)\|}_Z \le \|A\|\, {\|u+v\|}_{\gamma(U,X)},


$$


откуда получаем


$$


{\|Au\|}_Z \le \|A\| \inf \{ {\|u+v\|}_{\gamma(U,X)} : v \in


U_0 \} \le c {\|\gamma u\|}_X.


$$


Последнее неравенство справедливо в силу теоремы.


\end{pf}





\begin{examnonum}[{\series{m}\shape{n}\selectfont лемма Брэмбла-Гильберта}]


Применим следствие к


пространству Соболева (см. предыдущий пример). Пусть линейный


оператор $A$ непрерывно действует из $W^m_p(\Omega)$ в некоторое


нормированное пространство $Z$. Если ядро оператора $A$ содержит


пространство полиномов $P_m(\Omega)$, то имеет место неравенство


$$


{\|Au\|}_Z \le c \sum_{|i|=m} {\|D^i u\|}_{L_p(\Omega)}


\ \ \forall u \in W^m_p(\Omega).


$$


\end{examnonum}





\begin{thms} \label{renorm}


Пусть подпространство $U_0 =


\ker \gamma \cap U$ дополняемо в $V=\gamma(U,X)$ и $X \cap \gamma(U)$


замкнуто в $X$. Если непрерывная на $V$ полунорма $p$


эквивалентна на $U_0$ норме $U$, то функционал $u \to


p(u)+{\|\gamma u \|}_X$ является нормой на $V$, эквивалентной


норме усиленного пространства~$V$.


\end{thms}





{\bf Доказательство.} Функционал $u \to p(u)+{\|\gamma u \|}_X$,


очевидно, является нормой, которая в силу непрерывности


полунормы $p$ слабее нормы усиленного пространства $V$. Нужно


показать, что норма $V$ оценивается сверху нормой $u \to


p(u)+{\|\gamma u \|}_X$, помноженной на некоторую постоянную.


Пусть $\pi$ --- оператор проектирования $V$ на $U_0$. Тогда


оператор $A:V \to V$, действующий по формуле $Au=u - \pi u$,


непрерывен и удовлетворяет условию $Au=0$ при $u \in U_0$. В


силу следствия предыдущей теоремы имеем оценку


$$


{\|Au\|}_V \le c_0 {\|\gamma u\|}_X \ \ \forall u \in V.


$$


Используя эту оценку и то, что полунорма $p$ на $U_0$


эквивалентна норме $U$, получаем цепочку неравенств


\begin{multline*}


{\|u\|}_V \le {\|\pi u\|}_V+{\|Au\|}_V \le c_1p(\pi


u)+{\|Au\|}_V \le c_1p(u)+c_1p(Au)+{\|Au\|}_V \le \\


%


\le c_1p(u)+c_2{\|Au\|}_V \le c_1p(u)+c_2c_0{\|\gamma u\|}_X


\le c(p(u)+{\|\gamma u\|}_X).\quad\square


\end{multline*}





\begin{cornonum}


Пусть $\gamma(U,X)$


компактно вложено в $U$. Если непрерывная на $\gamma(U,X)$


полунорма~$p$ такова, что $\ker p \cap U_0 =\{0\}$, то функционал


$u \to p(u)+{\|\gamma u \|}_X$ является нормой, эквивалентной


норме усиленного пространства $\gamma(U,X)$. 


\end{cornonum}





{\bf Доказательство.} В силу следствия теоремы \ref{tcomp}


подпространство $U_0$ конечномерно, поэтому из условия $\ker p


\cap U_0=\{0\}$ следует, что сужение $p$ на $U_0$ является


нормой. Так как на конечномерном пространстве все нормы


эквивалентны, то на $U_0$ полунорма $p$ эквивалентна норме


пространства $U$. Далее, в силу теоремы \ref{DeniLions}


подмножество $X \cap \gamma(U)$ замкнуто в $X$. Таким образом,


выполнены все условия теоремы \ref{renorm}, применение которой к


рассматриваемому случаю доказывает утверждение.





\begin{examnonum}[{\series{m}\shape{n}\selectfont теорема 


Соболева об эквивалентных нормировках}]


В качестве иллюстрации снова рассмотрим


пространство Соболева $W^m_p(\Omega)$ как усиленное пространство


$\gamma(U,X)$, где $\gamma u=(D^i u)_{|i|=m}$, $U=L_p(\Omega)$,


$X=L_p(\Omega)^N$. В данной реализации доказанное выше


утверждение дает хорошо известную теорему Соболева о


перенормировках: если полунорма $p$ непрерывна на


$W^m_p(\Omega)$ и $\ker p \cap P_m(\Omega)=\{0\}$, то норма


$$


u \to p(u)+\sum_{|i|=m} {\|D^i u\|}_{L_p(\Omega)}


$$


эквивалентна норме пространства $W^m_p(\Omega)$.


\end{examnonum}








\begin{thebibliography}{30}





\bibitem{Anton}


Антонцев С.Н., Мейрманов А.М. {\em Математические модели движения


поверхностных и подземных вод}. -- Новосибирск: Изд-во НГУ,


1979. -- 79 с.





\bibitem{GlazPav1}


Глазырина Л.Л., Павлова М.Ф. {\em Разностная схема решения задачи совместного


движения грунтовых и поверхностных вод}


// Изв. вузов. Математика. -- 1984. --


\No\,9. -- C.\,72--75.





\bibitem{GlazPav}


Глазырина Л.Л., Павлова М.Ф. {\em О разрешимости одного нелинейного


эволюционного неравенства теории совместного движения поверхностных


и подземных вод} // Изв. вузов. Математика. -- 1997. -- \No\,4. -- C.\,20--31.





\bibitem{Daut}


Даутов Р.З., Карчевский М.М., Федотов Е.М. {\em Об одном методе решения


трехмерных стационарных задач типа сосредоточенной емкости} // Tез. докл.


Всесоюзн. конф. ``Актуальные проблемы вычислительной и прикладной 


математики'', Новосибирск, ВЦ СО АН СССР. -- 1987. -- C.\,68.





\bibitem{Djak96}


D'akonov E.G. {\em Optimization in solving elliptic problems}.


-- Boca Raton, 1996. 





\bibitem{DjakDAN}


Дьяконов Е.Г. {\em Новых подход к краевым условиям Дирихле,


основанный на использовании усиленных пространств Соболева}


// Докл. РАН. -- 1997. -- T.\,352.


-- \No\,5. -- C.\,590--594.





\bibitem{Djak}


Дьяконов Е.Г. {\em Усиленные прoстранства Соболева и некоторые новые типы


эллиптических краевых задач} // Дифференц. уравнения. -- 1997. -- T.\,33. --


\No\,4. -- C.\,532--539.





\bibitem{DjakIzVuz}


Дьяконов Е.Г. {\em Оценки $N$-поперечников в смысле Колмогорова 


для некоторых компактов в усиленных пространствах Соболева}


// Изв. вузов. Математика. -- 1997. -- \No\,4. -- C.\,32--50.





\bibitem{Trib}


Трибель Х. {\em Теория интерполяции. Функциональные


пространства. Дифференциальные операторы}. -- М.: Мир, 1980. --


664 с.





\bibitem{Aron}


Aronszain N., Gagliardo E. {\em Interpolation spaces and


interpolation methods} // Ann. Mat. Pura Appl.


-- 1965. -- V.\,68(4) -- P.\,51--117.





\bibitem{Necas}


Necas J. {\em Les Methodes Directes en Theorie des


Equations Elliptiques}. -- Masson, Paris/Academia, Pragua, 1967. -- 346 p.





\bibitem{Iosida}


Иосида К. {\em Функциональный анализ}. -- М.: Мир, 1967. -- 624 с.





\bibitem{Sobolev}


Соболев С.Л. {\em Избранные вопросы теории функциональных пространств


и обобщенных функций}. -- М.: Наука, 1989. -- 254 с.





\end{thebibliography}





\vskip 2cm





\post{\it Казанский государственный}{\it Поступила}


\post{\it университет}{19.11.2001}





\label{end}


\end{document}


